\section{Objectifs}
\label{sec:objectifs}

Pour juger nos résultats, nous nous appuierons surtout sur le F1 macro mesuré sur le jeu de test, car il donne le même poids à chaque classe, y compris aux plus rares. Nous préférons fixer la valeur à atteindre après les premiers essais avec les modèles classiques : elle reposera ainsi sur un score réel.

\begin{center}
\begin{tabularx}{\textwidth}{>{\raggedright\arraybackslash}p{6cm}Y}
  \toprule
  \textbf{Objectif} & \textbf{Mesure de réussite} \\
  \midrule
  Décrire le jeu de données retenu & Provenance, nombre d'images, répartition par classe, résolution et qualité documentées \\
  Comparer SVM, Random Forest et k-NN sur des descripteurs HOG & Accuracy, precision, recall, F1 et matrice de confusion, sur les mêmes données \\
  Entraîner un CNN de zéro & F1 macro comparé au meilleur modèle classique, sur le même jeu de test \\
  Analyser les classes difficiles & Paires de classes confondues et images mal classées expliquées \\
  Tester le modèle final sur de nouvelles images & Le programme affiche par exemple \og Image \textrightarrow{} Voiture (92~\%) \fg{}, où 92~\% est un score de confiance \\
  Démontrer le modèle dans le simulateur & Chaque classe prédite déclenche une décision, et les fautes graves sont comptées \\
  \bottomrule
\end{tabularx}
\end{center}
